\section{Quantum Walks}
\label{sec:prelim-qw}
% Rewritten 20 Sep 2026 (author's instruction): the former Secs. 2.3 and 2.4
% merged and rewritten from scratch with EQUAL treatment of the continuous-
% and discrete-time walks. The path-integral material and the "further use
% cases" subsection are gone (kept in draft.tex.bak-pre-qwrewrite); search,
% optimisation and Shor are Chapter 3. Labels kept because later chapters
% cite them: sec:prelim-qw, sec:ctqw, sec:prelim-dtqw, sec:dtqw-line,
% sec:dtqw-graphs, sec:dtqw-coins, sec:szegedy, sec:staggered,
% sec:lackadaisical, eq:ctqw, eq:shiftline, eq:walkprob, eq:walkrecursion,
% eq:uk, eq:dispersion, eq:konno, eq:graphcoin, eq:flipflop,
% eq:movingshift, eq:u2coin, eq:grovercoin, eq:fouriercoin, eq:szegedy.
% Convention, here and in Ch. 4: H = H_P (x) H_C, position first, coin
% second, so a discrete step is U = S (I (x) C).

A quantum walk is the unitary evolution of a quantum system --- the walker
--- whose position is indexed by the vertices of a graph $G = (V, E)$, and
which moves amplitude only along edges. The classical random walk it
generalises is a Markov chain on $V$; the quantum walk replaces the
stochastic matrix by a unitary, and the probability vector by a state
vector. That replacement can be made in two ways, and the two are the
subject of this section. The discrete-time walk of Aharonov, Davidovich and
Zagury~\cite{Aharonov_1993} advances in steps, each a fixed unitary; the
continuous-time walk of Farhi and Gutmann~\cite{Farhi_1997} evolves under a
time-independent Hamiltonian built from the graph. They are treated here on
an equal footing, because both are used later: the discrete-time walk is the
model of the learning algorithms of Chapter~\ref{ch:methodology} and of the
search and order-finding constructions of Chapter~\ref{ch:algorithms}, and
the continuous-time walk is the model in which the QAOA mixer and the
constant-time search of that chapter are formulated. Both are universal for
quantum computation~\cite{PhysRevLett.102.180501, ett2010universal}, and the
two are related as models~\cite{Childs2010, strauch2006connecting}; those
relations are made precise at the end of the section. Introductory
treatments are given by Kempe~\cite{Kempe_2003} and
Portugal~\cite{Portugal_2018}, and the reviews of
Venegas-Andraca~\cite{Elias_2012} and Kadian et al.~\cite{Kadian_2021} survey
the field.

\subsection{Two Quantisations of the Random Walk}
\label{sec:two-quantisations}
The classical random walk of Section~\ref{sec:prelim-crw} comes in the same
two forms, and each quantum walk is the quantisation of one of them. The
discrete-time chain updates a probability vector by a stochastic matrix,
$p_{t+1} = P\,p_t$; the continuous-time chain lets it evolve under the
generator $-\gamma L$, with $L = D - A$ the graph Laplacian, $\dot{p} =
-\gamma L\,p$, so that $p_t = e^{-\gamma L t} p_0$. Quantising means acting on
amplitudes rather than probabilities, and the requirement that the evolution
be unitary is what makes the two cases different.

In continuous time the requirement is met at once: $L$ (or $A$) is Hermitian,
so $e^{-i\gamma L t}$ is unitary, and the Schr\"odinger equation with $H =
\gamma L$ is a walk on $G$ with no further ingredient. In discrete time it is
not: a stochastic matrix is not unitary, and Meyer~\cite{Meyer_1996} showed
that on the line no translation-invariant unitary that moves amplitude only
between neighbouring sites exists other than a bare shift, up to a phase.
The escape is an internal degree of freedom --- the coin --- whose state
decides the direction of motion and which is not measured between steps.
That is the whole structural difference between the two models: the
continuous-time walk lives on $\operatorname{span}\{\ket{v}\}$ and is fixed by
the graph, the discrete-time walk lives on a larger space and carries a
choice, the coin, at every vertex. The rest of the section develops each in
turn, then places them side by side.

\subsection{The Continuous-Time Quantum Walk}
\label{sec:ctqw}
Let $A$ be the adjacency matrix of $G$, $D$ the diagonal matrix of degrees
and $L = D - A$ the Laplacian. On $\mathcal{H} = \operatorname{span}\{\ket{v}
: v \in V\}$, of dimension $N = |V|$, the walk is
\begin{equation}
	\ket{\psi(t)} = e^{-iHt} \ket{\psi(0)}, \qquad H = \gamma L \quad \text{or} \quad H = -\gamma A,
	\label{eq:ctqw}
\end{equation}
for a hopping rate $\gamma > 0$, and the position distribution is $p_t(v) =
|\braket{v|\psi(t)}|^2$. On a regular graph the two Hamiltonians differ by a
multiple of the identity and give identical probabilities; on an irregular
graph they do not, and the choice is part of the model. Since $H$ is
Hermitian the evolution is solved by its spectral decomposition, $H =
\sum_j \lambda_j \ket{\lambda_j}\bra{\lambda_j}$,
\begin{equation}
	\braket{v|\psi(t)} = \sum_j e^{-i\lambda_j t} \braket{v|\lambda_j}\braket{\lambda_j|\psi(0)},
	\label{eq:ctqw-spectral}
\end{equation}
so every property of the walk is a property of the spectrum of the graph:
the walker's amplitude at $v$ is a superposition of oscillations at the
eigenfrequencies of $G$, weighted by the overlaps of the eigenvectors with
the start and end vertices. The classical chain has the same
decomposition with $e^{-\gamma\lambda_j t}$ in place of $e^{-i\lambda_j t}$:
the classical walk relaxes to the stationary distribution at a rate set by
the spectral gap, the quantum walk never relaxes, and its time-averaged
distribution, $\bar{p}(v) = \lim_{T\to\infty} T^{-1}\int_0^T p_t(v)\,dt$, is
what plays the role of a limit~\cite{Aharonov_2000}. The evolution commutes
with every automorphism of $G$, so the walk inherits all of the graph's
symmetry; there is no coin to break it.

\paragraph*{The walk on the line.}
On the infinite line, $A\ket{x} = \ket{x-1} + \ket{x+1}$, the momentum states
$\ket{k} = \sum_x e^{ikx}\ket{x}$, $k \in [-\pi, \pi)$, diagonalise $A$ with
$A\ket{k} = 2\cos k\,\ket{k}$, so with $H = -\gamma A$ the dispersion relation
is
\begin{equation}
	E(k) = -2\gamma \cos k, \qquad v_g(k) = \frac{dE}{dk} = 2\gamma \sin k,
	\label{eq:ctqw-dispersion}
\end{equation}
and the amplitude of a walker started at the origin is a Bessel function,
\begin{equation}
	\braket{x|\psi(t)} = \frac{1}{2\pi}\int_{-\pi}^{\pi} e^{ikx + 2i\gamma t\cos k}\,dk = i^{x} J_x(2\gamma t),
	\qquad p_t(x) = J_x(2\gamma t)^2 .
	\label{eq:ctqw-bessel}
\end{equation}
The group velocity is bounded by $2\gamma$, and the stationary-phase
analysis of equation~\eqref{eq:ctqw-bessel} places the probability mass in
$|x| \le 2\gamma t$ with sharp fronts at the edges: the spread is
\emph{ballistic}, the standard deviation of $x$ growing as $\sqrt{2}\gamma t$
rather than as $\sqrt{t}$. Konno~\cite{Konno_2005} proved the corresponding
weak limit: $X_t/t$ converges in distribution to the arcsine density
\begin{equation}
	f_C(x) = \frac{1}{\pi\sqrt{4\gamma^2 - x^2}}, \qquad |x| < 2\gamma ,
	\label{eq:ctqw-limit}
\end{equation}
bimodal, diverging at the fronts $\pm 2\gamma$, and --- because there is no
coin --- the same for every initial state localised at a point. This is the
continuous-time counterpart of the Konno law~\eqref{eq:konno} below, and the
comparison is instructive: both walks spread linearly in time and both limit
densities blow up at the fronts, but the discrete-time density depends on
the coin and the initial coin state and the continuous-time one does not.

\paragraph*{Where it has been decisive.}
The continuous-time walk produced the strongest separation known between a
walk and any classical algorithm, the exponential speed-up of the
glued-trees traversal of Childs et al.~\cite{Childs_2001, Childs_2002}, and
the first universality proof~\cite{PhysRevLett.102.180501}, later extended to
multi-particle walks~\cite{Childs_2012}. It is the natural language for
coherent transport on networks~\cite{Mlken_2011}, and, because the
transverse-field mixer of the quantum approximate optimisation algorithm is
$e^{-i\beta A}$ for the hypercube, it is the model in which QAOA is a walk
(Section~\ref{sec:optimisation}).

\subsection{The Discrete-Time Quantum Walk}
\label{sec:prelim-dtqw}
The discrete-time walk carries a coin. The convention throughout the thesis
is that the Hilbert space is the tensor product of a position space and a
coin space, in that order,
\begin{equation}
	\mathcal{H} = \mathcal{H}_P \otimes \mathcal{H}_C ,
\end{equation}
so that one step is $U = S\,(I \otimes C)$: the coin operator $C$ acts on the
second factor and the shift $S$ on both. On the line, $\mathcal{H}_P =
\operatorname{span}\{\ket{x} : x \in \mathbb{Z}\}$ and $\mathcal{H}_C =
\mathbb{C}^2$ with basis $\{\ket{0}, \ket{1}\}$ read as ``move right'' and
``move left''. The coin is any $C \in U(2)$, the canonical choice being the
Hadamard coin $H = \tfrac{1}{\sqrt{2}}\bigl(\begin{smallmatrix} 1 & 1 \\ 1 &
-1 \end{smallmatrix}\bigr)$, and the shift moves the walker in the direction
the coin indicates,
\begin{equation}
	S = \sum_{x \in \mathbb{Z}} \Bigl( \ket{x+1}\bra{x} \otimes \ket{0}\bra{0} + \ket{x-1}\bra{x} \otimes \ket{1}\bra{1} \Bigr).
	\label{eq:shiftline}
\end{equation}
After $t$ steps from $\ket{\psi_0}$ the state is $U^t\ket{\psi_0}$ and the
position distribution is obtained by tracing out the coin,
\begin{equation}
	p_t(x) = \sum_{c \in \{0, 1\}} \bigl| \braket{x, c | \psi_t} \bigr|^2 ;
	\label{eq:walkprob}
\end{equation}
in components $\psi_t(x, c) = \braket{x, c|\psi_t}$ one step is the pair of
recursions
\begin{equation}
	\psi_{t+1}(x, 0) = \sum_{c} C_{0c}\, \psi_t(x-1, c), \qquad
	\psi_{t+1}(x, 1) = \sum_{c} C_{1c}\, \psi_t(x+1, c).
	\label{eq:walkrecursion}
\end{equation}
The coin is tossed by $C$ and the walker moves conditioned on the outcome,
but the outcome is not measured: if it were, the state would collapse to a
definite position at each step and equation~\eqref{eq:walkrecursion} would
reduce, for the Hadamard coin, to the classical symmetric walk $p_{t+1}(x) =
\tfrac12 p_t(x-1) + \tfrac12 p_t(x+1)$~\cite{Kempe_2003}. The quantum walk is
what remains when the measurements are omitted and the amplitudes interfere.
Interference is visible already at three steps. From $\ket{0}\otimes\ket{0}$
with the Hadamard coin,
\begin{equation}
	\ket{\psi_3} = \frac{1}{2\sqrt{2}} \Bigl( \ket{3, 0} + \ket{1, 1} + 2\ket{1, 0} - \ket{-1, 0} + \ket{-3, 1} \Bigr),
\end{equation}
so $p_3(1) = \tfrac58$ while $p_3(\pm 3) = p_3(-1) = \tfrac18$: the two
routes to $x = 1$ add, those to $x = -1$ partly cancel, and the distribution
is skewed to the right. Starting from $\ket{0}\otimes\ket{1}$ gives the
mirror image, and the coin state $(\ket{0} + i\ket{1})/\sqrt{2}$ a symmetric
distribution~\cite{Kempe_2003}. The dependence of the outcome on the coin
state, absent in both the classical walk and the continuous-time quantum
walk, is the first indication that the coin is a substantive parameter of
the model~\cite{Tregenna_2003}.

\subsection{Exact Solution on the Line and the Limit Laws}
\label{sec:dtqw-line}
Because $U$ commutes with translations it is diagonalised in the position
factor by the same Fourier transform as the continuous-time walk. With
$\ket{k} = \sum_x e^{ikx}\ket{x}$, the shift~\eqref{eq:shiftline} acts as
$S(\ket{k}\otimes\ket{0}) = e^{-ik}\ket{k}\otimes\ket{0}$ and
$S(\ket{k}\otimes\ket{1}) = e^{ik}\ket{k}\otimes\ket{1}$, so on each momentum
sector the step is the $2 \times 2$ unitary
\begin{equation}
	U_k = \begin{pmatrix} e^{-ik} & 0 \\ 0 & e^{ik} \end{pmatrix} C ,
	\qquad
	U_k^{(H)} = \frac{1}{\sqrt{2}} \begin{pmatrix} e^{-ik} & e^{-ik} \\ e^{ik} & -e^{ik} \end{pmatrix}
	\label{eq:uk}
\end{equation}
for the Hadamard coin, whose eigenvalues are $e^{-i\omega_k}$ and
$-e^{i\omega_k}$ with
\begin{equation}
	\sin \omega_k = \frac{\sin k}{\sqrt{2}}, \qquad \omega_k \in [-\tfrac{\pi}{2}, \tfrac{\pi}{2}],
	\label{eq:dispersion}
\end{equation}
as one checks from the trace $-i\sqrt{2}\sin k$ and determinant $-1$
of~\eqref{eq:uk}. Equation~\eqref{eq:dispersion} is the discrete-time
dispersion relation, the counterpart of~\eqref{eq:ctqw-dispersion}, and the
analysis of Ambainis et al.~\cite{Ambainis_2001} and Nayak and
Vishwanath~\cite{Nayak_2000} proceeds from it by stationary phase exactly as
for the Bessel integral~\eqref{eq:ctqw-bessel}: the amplitude at $x$ after
$t$ steps is an integral over $k$ of terms $e^{i(kx \mp \omega_k t)}$, whose
stationary points move at the group velocity $d\omega_k/dk = \cos k/\sqrt{2 -
\sin^2 k}$, maximal at $1/\sqrt{2}$. The probability mass therefore lies, up
to exponentially small corrections, in $|x| \le t/\sqrt{2}$, nearly flat
inside with sharp peaks at the two ends, and the standard deviation grows
linearly in $t$.

The precise statement is Konno's weak limit theorem~\cite{Konno_2002},
generalised by Grimmett, Janson and Scudo~\cite{Grimmett_2003}: for the
Hadamard walk $X_t/t$ converges in distribution, and for the symmetric
initial coin state the limiting density is
\begin{equation}
	f_K(x) = \frac{1}{\pi (1 - x^2) \sqrt{1 - 2x^2}}, \qquad |x| < \frac{1}{\sqrt{2}},
	\label{eq:konno}
\end{equation}
bimodal and diverging at $\pm 1/\sqrt{2}$. For a general initial coin state
the density is multiplied by $1 - \kappa x$ with $|\kappa| \le 1$ fixed by
that state, which is the asymptotic form of the skew seen at three steps.
Set beside~\eqref{eq:ctqw-limit}, the two limit laws say the same thing
about speed --- both walks are ballistic, both the classical $\sqrt{t}$ ---
and different things about shape: the continuous-time density is the
arcsine law and depends on nothing but $\gamma$, the discrete-time density
has the extra factor $(1 - x^2)^{-1}$ and the coin-dependent skew. Knight,
Rold\'an and Sipe~\cite{Knight_2003} showed that in both cases the phenomenon
is wave interference and requires no quantisation of the walker as such.
Two consequences matter later. Ballistic spreading is what lets a walk of
$t$ steps ``see'' a neighbourhood of radius proportional to $t$, and is the
origin of the quadratic speed-ups of Section~\ref{sec:search}; and the
sensitivity of $\kappa$, and of the whole dispersion
relation~\eqref{eq:dispersion}, to the coin is what makes the coin a design
variable in the discrete-time model, with no analogue in the continuous-time
one.

\subsection{Walks on General Graphs}
\label{sec:dtqw-graphs}
The continuous-time walk needs nothing further on a general graph:
equation~\eqref{eq:ctqw} is defined for any $G$, and
equation~\eqref{eq:ctqw-spectral} solves it. The discrete-time walk does,
because an arbitrary vertex has no canonical ``left'' and ``right''. It is
most cleanly defined on the space of directed
arcs~\cite{Aharonov_2000, Kempe_2003, Portugal_2018}: let $\mathcal{H} =
\operatorname{span}\{\ket{v, w} : v \in V,\ w \in N(v)\}$, of dimension
$2|E|$ for a simple graph, where $\ket{v, w}$ is read ``at $v$, pointing
towards $w$''. The subspace $\mathcal{H}_v = \operatorname{span}\{\ket{v, w}
: w \in N(v)\}$ has dimension $\deg v$ and is the coin space at $v$, so that
$\mathcal{H} = \bigoplus_v \mathcal{H}_v$ is what $\mathcal{H}_P \otimes
\mathcal{H}_C$ becomes when the degree varies. The coin is a direct sum of
local unitaries,
\begin{equation}
	C = \bigoplus_{v \in V} C_v, \qquad C_v \in U(\deg v),
	\label{eq:graphcoin}
\end{equation}
each redistributing the amplitude at $v$ among its outgoing arcs, and two
shifts are in use. The flip-flop shift reverses the arc,
\begin{equation}
	S_{\mathrm{ff}} \ket{v, w} = \ket{w, v} ;
	\label{eq:flipflop}
\end{equation}
it is defined on any graph, is an involution, and is a permutation of the
arc basis, so it is free in the sense of
Section~\ref{sec:efficientcircuit}. The moving shift requires a $d$-regular
graph whose arcs are coloured so that the arc $(v, w)$ and the arc leaving
$w$ ``in the same direction'' share a colour --- on the line the colours are
right and left, on the hypercube the bit flipped --- and then, with a single
coin space $\mathbb{C}^d$ shared by all vertices,
\begin{equation}
	S_{\mathrm{m}} \ket{v} \otimes \ket{c} = \ket{w_c(v)} \otimes \ket{c},
	\label{eq:movingshift}
\end{equation}
where $w_c(v)$ is the $c$-coloured neighbour of $v$. Equation~\eqref{eq:shiftline}
is the moving shift on the line, and the colour-by-colour decomposition of
equation~\eqref{eq:colourshift} in Chapter~\ref{ch:methodology} is the same
operator written for a circuit. The two shifts are not interchangeable:
Ambainis, Kempe and Rivosh~\cite{Ambainis_2004} found that search on the
two-dimensional grid with the Grover coin achieves its $O(\sqrt{N\log N})$
running time only with the flip-flop shift and gives no speed-up with the
moving one. A step of the general walk is $U = S\,C$ with $C$
from~\eqref{eq:graphcoin} and $S$ from~\eqref{eq:flipflop}
or~\eqref{eq:movingshift}, and the position distribution is $p_t(v) =
\sum_{w \in N(v)} |\braket{v, w|\psi_t}|^2$.

The contrast with the continuous-time walk on the same graph is now exact.
Both walks are local: the continuous-time walk because $H$ has no matrix
elements between non-adjacent vertices, the discrete-time walk because each
step moves amplitude along one arc. Both are diagonalised by the graph's
symmetries where the graph has them. But the continuous-time walk is
determined by $G$ and one number, $\gamma$, while the discrete-time walk on
the same $G$ carries $\sum_v \dim U(\deg v)$ real parameters in its coins and
a further choice of shift; the walk that Chapter~\ref{ch:methodology} treats
as a design space is the discrete-time one for exactly this reason.

\subsection{Coin Operators}
\label{sec:dtqw-coins}
Three coin families dominate the literature, and the methodology measures
its own coins against them. For $d = 2$ every coin is, up to a global phase,
a member of the three-parameter family
\begin{equation}
	C(\theta, \phi, \lambda) = \begin{pmatrix} \cos\theta & e^{i\phi} \sin\theta \\ e^{i\lambda} \sin\theta & -e^{i(\phi + \lambda)} \cos\theta \end{pmatrix},
	\label{eq:u2coin}
\end{equation}
of which the Hadamard coin is $\theta = \pi/4$, $\phi = \lambda = 0$; $\theta$
controls the balance between continuing and reversing, the phases the
interference between the two, and the average position and spread of the
walk have been mapped across the whole family~\cite{Li_2012, Souza_2013}.
For $d > 2$ the standard choice is the Grover coin, the reflection about the
uniform coin state $\ket{s} = d^{-1/2}\sum_c \ket{c}$,
\begin{equation}
	G = 2\ket{s}\bra{s} - I, \qquad G_{cc'} = \frac{2}{d} - \delta_{cc'} ,
	\label{eq:grovercoin}
\end{equation}
real, symmetric, invariant under every permutation of the coin states and
the permutation-symmetric real coin farthest from the identity; it is the
coin of the search algorithms of Section~\ref{sec:search}, and for $d = 2$ it
degenerates to Pauli $X$. The discrete Fourier transform on $\mathbb{C}^d$,
\begin{equation}
	F_{cc'} = \frac{1}{\sqrt{d}}\, \omega^{c c'}, \qquad \omega = e^{2\pi i / d},
	\label{eq:fouriercoin}
\end{equation}
is the standard complex coin; for $d = 2$ it coincides with Hadamard, and its
spectrum differs qualitatively from the Grover coin's~\cite{Komatsu_2019,
Saito_2018}. That the choice matters is settled~\cite{Tregenna_2003}; what it
changes in a learning model is the question of Chapter~\ref{ch:methodology},
whose coin $K(\delta, \theta)$ is defined in Section~\ref{sec:coindesign}.

\subsection{Relations Between the Models, and Alternative Formulations}
\label{sec:dtqw-alternatives}
The two walks are not independent models. Strauch~\cite{strauch2006connecting}
showed that the continuous-time walk on the line is a limit of the
discrete-time walk: with a coin close to the identity, $C = e^{-i\epsilon
\sigma}$ for a suitable $\sigma$, and time rescaled as $t = \tau/\epsilon$,
the discrete-time dynamics converge to Hamiltonian evolution as $\epsilon \to
0$. Childs~\cite{Childs2010} showed the converse direction in general: any
continuous-time walk $e^{-iHt}$ can be simulated by a discrete-time walk,
namely Szegedy's walk built from a Markov chain derived from $H$, with the
eigenphases of the discrete walk related to the eigenvalues of $H$ by an
$\arcsin$; this is how Hamiltonian simulation by quantum walk works, and it
is the sense in which the committee's remark that the two models are
``functionally equivalent'' is correct. It is not, however, an equivalence of
cost: the continuous-time walk is Hamiltonian evolution and running it on a
gate-model device requires Hamiltonian simulation --- circuit constructions
exist for sparse~\cite{Chen_2024a} and random~\cite{Chakraborty_2025} graphs,
but they approximate the walk --- whereas the discrete-time walk is already a
sequence of unitaries and compiles to gates exactly~\cite{Douglas_2007,
Acasiete_2020}. That asymmetry is why the simulations of
Chapters~\ref{ch:methodology} and~\ref{ch:results} are built from the
discrete-time walk, and why the continuous-time walk is used in
Chapter~\ref{ch:algorithms} only where the construction in the literature
is continuous-time.

Three reformulations of the discrete-time walk recur in what follows.

\paragraph*{Szegedy's quantisation of a Markov chain.}
\label{sec:szegedy}
Szegedy~\cite{Szegedy} builds a walk from a stochastic matrix $P$ on $N$
states rather than from a graph. On $\mathbb{C}^N \otimes \mathbb{C}^N$
define $\ket{\psi_x} = \ket{x} \otimes \sum_y \sqrt{P_{xy}}\,\ket{y}$ and
$\ket{\phi_y} = \sum_x \sqrt{P_{yx}}\,\ket{x} \otimes \ket{y}$, the projectors
$\Pi_A = \sum_x \ket{\psi_x}\bra{\psi_x}$ and $\Pi_B = \sum_y
\ket{\phi_y}\bra{\phi_y}$, and the walk operator
\begin{equation}
	W = (2\Pi_B - I)(2\Pi_A - I),
	\label{eq:szegedy}
\end{equation}
a product of two reflections. If $\cos\theta_j$ are the singular values of
the discriminant $D_{xy} = \sqrt{P_{xy}P_{yx}}$, the eigenvalues of $W$ on
the span of the $\ket{\psi_x}$ and $\ket{\phi_y}$ are $e^{\pm 2i\theta_j}$:
a spectral gap $\delta$ of a symmetric $P$ becomes a phase gap of order
$\sqrt{\delta}$, which is the source of the quadratic speed-ups in hitting
and detection~\cite{Magniez_2007}. Wong~\cite{G_2017} showed the construction
equivalent to the coined walk with the Grover coin~\eqref{eq:grovercoin} and
the flip-flop shift~\eqref{eq:flipflop} on the bipartite double cover, so the
coined walk of Section~\ref{sec:dtqw-graphs} includes Szegedy's; it is also
the discrete-time walk that simulates a continuous-time one in Childs's
construction.

\paragraph*{The staggered model.}
\label{sec:staggered}
Portugal et al.~\cite{Portugal_2016} dispense with the coin and act on the
position space alone. A tessellation of $G$ is a partition of $V$ into
cliques, and a tessellation cover $\mathcal{T}_1, \ldots, \mathcal{T}_m$ is a
set of them such that every edge lies inside a clique of at least one. With
$\ket{\alpha} = |\alpha|^{-1/2}\sum_{v\in\alpha}\ket{v}$ for each clique, the
reflections $U_k = 2\sum_{\alpha \in \mathcal{T}_k}\ket{\alpha}\bra{\alpha} -
I$ each move amplitude only within their cliques, and $U = U_m \cdots U_1$ is
one step. Szegedy's walk is the two-tessellation case on the double cover,
and the coinless lattice walks of Patel et al.~\cite{Patel_2004} are
staggered walks. The model shows that the coin is not essential to having a
discrete-time walk, only to having a homogeneous one, and that ``coin'' and
``tessellation'' are two ways of supplying the same local unitary freedom.

\paragraph*{The lackadaisical walk.}
\label{sec:lackadaisical}
Wong~\cite{Wong_2015, G_2018} adds to every vertex a self-loop of weight $\ell
> 0$. The coin space at $v$ acquires the arc $\ket{v, v}$, the Grover coin is
taken about the weighted uniform state
\begin{equation}
	\ket{s_v} = \frac{1}{\sqrt{\deg v + \ell}} \Bigl( \sum_{w \in N(v)} \ket{v, w} + \sqrt{\ell}\, \ket{v, v} \Bigr), \qquad
	C_v = 2\ket{s_v}\bra{s_v} - I,
	\label{eq:lackadaisical}
\end{equation}
and the flip-flop shift fixes the self-loop arc. The single scalar $\ell$
decides whether search works: on the two-dimensional grid $\ell = 4/N$
raises the success probability from $O(1/\log N)$ to nearly
one~\cite{G_2018}, and the literature on which weights work on which
graphs~\cite{Giri_2020, Nahimovs_2018, Rhodes_2019, Rapoza_2021} is the
clearest demonstration in the search setting that a parameter inside the
coin decides whether an algorithm works. It is the direct precedent for
treating coin parameters as tunable in a learning setting, and it has no
continuous-time counterpart other than $\gamma$ itself.

